% ============================================================
% tables_deployment_audit.tex
% Deployment measurement + early-exit audit tables.
% Label prefix: tab:dep:*    Macro prefix: \DEP*
% Requires: booktabs, pifont (\ding{51}/\ding{55})
%
% Source artifacts:
%   results/deployment/deployment_summary.csv
%   results/deployment/energy_model.csv
%   results/deployment/energy_crossover.csv
%   results/evidence_arrival.json
%   results/early_exit_audit.log
%
% NOT part of the original Table III/VI/VII/VIII series.
% Artifact sizes are decimal (bytes/1000), means over the eleven benchmarks.
% ============================================================

\begin{table}[!t]
\centering
\caption{Measured deployment characteristics of seven candidate detectors.
Macro-F1 is the mean over the eleven benchmarks; artifact size is the serialized
on-disk footprint of every component required at inference, measured from real
files rather than inferred from parameter counts; latency is median wall-clock
per sample on the same CPU. Operation counts are per sample: \emph{cmp} denotes a
tree node comparison, \emph{MAC} a dense multiply--accumulate, and \emph{SOP} a
spike-triggered accumulate. Arithmetic energy prices every operation under the
Horowitz 45\,nm model used throughout this paper, with a tree comparison charged
at the accumulate cost, since a subtract-and-branch belongs to that cost class.
Priced this way the compact ensemble requires roughly half the arithmetic energy
of the deployed spiking student.}
\label{tab:dep:systems}
\small
\begin{tabular}{lrrrlrr}
\toprule
 & F1-M & Artifact & $p_{50}$ & Operations & \multicolumn{2}{c}{$E_\mathrm{arith}$ (pJ)} \\
\cmidrule(lr){6-7}
System & (\%) & & (ms) & (per sample) & fp32 & int8 \\
\midrule
full\_gbdt     & 82.13 & 5.04\,MB  & 2.332  & 7,636\,cmp  & 6,873  & 229   \\
compact\_gbdt  & 81.32 & 2.41\,MB  & 2.187  & 3,076\,cmp  & 2,768  & 92    \\
dnn\_tree      & 79.86 & 50.2\,KB  & 0.226  & 8,471\,MAC  & 39,140 & 1,954 \\
t2s\_full      & 79.80 & 53.4\,KB  & 2.909  & 7,587\,SOP  & 7,000  & 233   \\
t2s\_exit      & 79.68 & 53.4\,KB  & 3.199  & 6,091\,SOP  & 5,654  & 189   \\
raw\_mlp\_fp   & 77.20 & 47.6\,KB  & 0.0135 & 10,979\,MAC & 50,503 & 2,525 \\
raw\_mlp\_int8 & 76.23 & 47.6\,KB  & 0.0135 & 10,979\,MAC & 50,503 & 2,525 \\
\bottomrule
\end{tabular}
\end{table}

\begin{table}[!t]
\centering
\caption{Sensitivity of the Tree-to-Spike versus DNN-Tree total-energy ratio to
the memory-access cost and the cache model. Values above unity favour
Tree-to-Spike. The cold model charges one access per operation with no reuse; the
reuse-aware model charges the first touch of a distinct item at the tier matching
the artifact size and subsequent touches at an L1 cost of 1\,pJ. Memory accounts
for 95.7--99.9\% of total energy in every configuration, so the ordering is
determined by the memory model rather than by operation counts. Under the cold
model the ordering \emph{inverts} above roughly 20\,pJ per access; under the
reuse-aware model it does not invert anywhere in the swept range. The 53\,KB
working set exceeds typical L1 capacity but fits L2; on targets without L2,
``reuse'' denotes SRAM residency rather than caching. The reported band
0.94--2.06 therefore spans a modelling choice, not measurement noise.}
\label{tab:dep:sensitivity}
\small
\begin{tabular}{lrrrrr}
\toprule
 & \multicolumn{5}{c}{SRAM cost (pJ per access)} \\
\cmidrule(lr){2-6}
Memory model & 5 & 10 & 20 & 50 & 100 \\
\midrule
Cold (1 access/op) & 1.54 & 1.24 & 1.08 & 0.98 & 0.94 \\
Reuse-aware        & 2.06 & 1.77 & 1.59 & 1.48 & 1.43 \\
\bottomrule
\end{tabular}
\end{table}

\begin{table}[!t]
\centering
\caption{Early-exit audit across the eleven benchmarks. \emph{Gate} indicates
whether the validation criterion admitted early exit. \textbf{$t'$ is the
evidence-arrival step, defined as the first timestep at which the correct class
becomes the running argmax under the full window, computed only over samples that
are correctly classified at $T$}; Incl.\ reports the share of test samples meeting
that condition, and equals full-window accuracy. This conditioning matters: on
NSL-KDD and UNSW-NB15 the samples excluded from $t'$ are largely the same samples
that exit incorrectly, so a low $t'$ beside a high Wrong\% is not a contradiction.
Undef.\ counts classes for which fewer than 5\% of samples are ever correctly
classified, which reflects a taxonomy too fine for the sample budget rather than
an early-exit failure. Wrong\% is the fraction of early exits disagreeing with the
full-window prediction, and the modal target is the class those exits most often
land on. The gate tracks wrong-exit \emph{volume} rather than concentration:
CICIoMT2024 passes despite 93.1\% concentration because only 2.1\% of exits are
wrong, whereas ISCXIDS2012 fails at 1.0\% because macro-F1 on a strongly
imbalanced binary task is sensitive enough that this breaches the one-point bound.
Because the $t'$ distributions are strongly right-skewed, the median and $p_{90}$
are the operationally meaningful statistics; the mean is reported for
completeness only.}
\label{tab:dep:exitaudit}
\small
\begin{tabular}{llrrrrrrlr}
\toprule
 & & \multicolumn{3}{c}{$t'$} & Incl. & Undef. & Wrong & Modal & Share \\
\cmidrule(lr){3-5}
Dataset & Gate & mean & med & $p_{90}$ & (\%) & cls. & (\%) & target & (\%) \\
\midrule
CICIoT2023      & \ding{55} & 4.83 & 1.0 & 13.0 & 98.8 & 0 & 44.8 & DoS        & 99.3 \\
NSL-KDD         & \ding{55} & 1.27 & 1.0 &  1.0 & 76.8 & 1 & 20.7 & normal     & 85.4 \\
UNSW-NB15       & \ding{55} & 4.87 & 3.0 & 15.0 & 73.7 & 3 & 19.2 & Normal     & 52.9 \\
Edge-IIoTset    & \ding{55} & 3.46 & 1.0 & 14.0 & 98.0 & 0 & 12.7 & Ransomware & 50.5 \\
ISCXIDS2012     & \ding{55} & 1.12 & 1.0 &  1.0 & 99.7 & 0 &  1.0 & attack     & 98.2 \\
\midrule
CIC-IDS2017     & \ding{51} & 2.81 & 2.0 &  6.0 & 99.8 & 0 &  5.5 & dos        & 86.8 \\
KDDCup99        & \ding{51} & 1.04 & 1.0 &  1.0 & 93.0 & 0 &  6.9 & normal     & 91.4 \\
CSE-CIC-IDS2018 & \ding{51} & 2.02 & 1.0 &  4.0 & 98.1 & 0 &  0.2 & normal     & 97.3 \\
TON-IoT         & \ding{51} & 5.09 & 5.0 & 12.0 & 96.8 & 0 &  3.9 & password   & 24.0 \\
CICIoMT2024     & \ding{51} & 1.43 & 1.0 &  2.0 & 99.7 & 0 &  2.1 & DoS        & 93.1 \\
CTU-13          & \ding{51} & 1.01 & 1.0 &  1.0 &100.0 & 0 &  0.0 & normal     & 59.8 \\
\bottomrule
\end{tabular}
\end{table}

\begin{table}[!t]
\centering
\caption{Per-class exit behaviour on CICIoT2023 at $\theta_\mathrm{exit}=0.90$:
the sibling-attractor failure. Recall under early exit collapses for DDoS while
the two classes it is mistaken for are unaffected, and 99.3\% of incorrect exits
land on DoS. Of samples that exit incorrectly, 99.3\% eventually reach the correct
class if the window is allowed to run, at a mean recovery step of 11.16 against a
mean wrong-exit step of 2.25 --- the evidence is present but arrives after the
gate has committed.}
\label{tab:dep:ciciot_perclass}
\small
\begin{tabular}{lrrrrr}
\toprule
Class & $n$ & Exit (\%) & $\bar{t}^*$ & Rec.\ exit & Rec.\ full \\
\midrule
DDoS       & 339,727 &  99.9 &  6.22 & 38.95 & 99.58 \\
DoS        &  80,877 & 100.0 &  4.15 & 99.95 & 99.64 \\
Mirai      &  26,427 &  99.6 &  3.24 & 99.27 & 99.61 \\
Benign     &  10,981 &  94.0 &  9.90 & 99.13 & 99.57 \\
Spoofing   &   4,880 &  67.3 & 11.46 & 70.64 & 70.68 \\
Recon      &   3,604 &  63.8 & 12.53 & 46.73 & 51.78 \\
Web        &     245 &  24.5 & 17.82 &  7.35 &  7.35 \\
BruteForce &     125 &  38.4 & 16.63 & 16.80 & 16.80 \\
\bottomrule
\end{tabular}
\end{table}

\begin{table}[!t]
\centering
\caption{Per-class exit behaviour on Edge-IIoTset at $\theta_\mathrm{exit}=0.90$:
a distinct, diffuse failure. Unlike CICIoT2023, no single class absorbs the
incorrect exits (modal share 50.5\%); instead recall collapses independently
across five classes, three of them near-total. This is late separation spread over
a fine taxonomy rather than commitment to one confusable sibling, and it is why a
single mechanism does not explain every disabled dataset.}
\label{tab:dep:edgeiiot_perclass}
\small
\begin{tabular}{lrrrrr}
\toprule
Class & $n$ & Exit (\%) & $\bar{t}^*$ & Rec.\ exit & Rec.\ full \\
\midrule
Normal         & 272,800 & 100.0 &  6.17 & 100.00 & 100.00 \\
DDoS\_UDP      &  24,314 & 100.0 &  9.58 &  97.77 &  99.95 \\
DDoS\_ICMP     &  13,588 & 100.0 &  2.59 &   5.69 &  99.52 \\
DDoS\_TCP      &  10,012 & 100.0 &  5.09 &   0.00 & 100.00 \\
Password       &   9,987 &  85.9 &  7.10 &  25.56 &  82.04 \\
DDoS\_HTTP     &   9,709 &  54.4 & 15.11 &  95.32 &  95.74 \\
Backdoor       &   4,805 & 100.0 &  1.46 &   1.29 &  93.74 \\
Port\_Scanning &   3,995 & 100.0 &  4.62 &  49.69 & 100.00 \\
Fingerprinting &     171 &  97.7 &  2.58 &   2.34 &  64.91 \\
MITM           &      72 &  90.3 &  3.75 &  91.67 &  91.67 \\
\bottomrule
\end{tabular}
\end{table}

\begin{table}[!t]
\centering
\caption{Exit-statistic exhaustion on CICIoT2023 (D2). If the failure were that
the max-posterior gate ignores a close runner-up, the margin
$q^{(t)}_{(1)}-q^{(t)}_{(2)}$ would be smaller at incorrect exits. It is not: the
separation is $-0.006$, in the \emph{wrong} direction, and a margin gate breaches
the one-point bound at every threshold swept. The prefix posterior is therefore
not an uncertain estimate of the full-window answer; it is a confident estimate of
a different quantity, and no statistic computed from the same prefix can separate
the two cases. This is the diagnostic signature of target mismatch as opposed to
incompleteness: adding information moves the estimate, whereas exhausting the
statistic space over fixed information does not.}
\label{tab:dep:margin}
\small
\begin{tabular}{lrr}
\toprule
Margin at exit & mean & median \\
\midrule
Correct exits   & 0.8730 & 0.8663 \\
Incorrect exits & 0.8791 & 0.8727 \\
Separation      & $-0.0061$ & $-0.0064$ \\
\midrule
\multicolumn{3}{l}{\emph{Margin gate }$q_{(1)}-q_{(2)}\ge m$} \\
$m$ & mean exit & F1-M gap \\
\midrule
0.50 &  2.12 & $-39.97$ \\
0.90 &  7.20 & $-36.60$ \\
0.99 & 11.21 & $-24.53$ \\
\bottomrule
\end{tabular}
\end{table}
